\section{Testing, coverage, and engineering infrastructure}\label{sec:testing}

The engineering infrastructure is the part of QtRocket that least resembles a young project:
a single CMake project pinned to C++23 (\src{CMakeLists.txt}{5}), every dependency except
Qt6 vendored at an exact tag, warnings-as-errors on four compiler families, 247 test cases
checked largely against independent oracles rather than golden values, a five-platform CI
matrix, and branch-level source coverage wired into the build. The same inspection finds the
document's recurring motif in infrastructural form --- machinery finished to a high standard
with no consumer: the coverage target is run by no CI job, the translation rules produce
files no target references, the install rule ships one of three executables, and the project
version surfaces nowhere a user could see it. This section reviews both halves.

\subsection{Build system: presets, pinned dependencies, strict warnings}\label{sec:testing-build}

\srcf{CMakePresets.json} defines eight configure presets --- \code{debug}, \code{release},
\code{debug-clang}, \code{release-clang}, \code{debug-gcc}, \code{release-gcc},
\code{release-msvc}, and \code{coverage-clang} --- every one pinned to the Ninja generator
(\src{CMakePresets.json}{7}) and paired with matching build and test presets. MSVC exists
only as a release variant (\src{CMakePresets.json}{70}); there is no \code{debug-msvc}. The
coverage preset inherits \code{debug-clang} but redirects to a separate
\code{build-coverage/} tree (\src{CMakePresets.json}{35}) and sets
\code{QTROCKET\_ENABLE\_COVERAGE} (\src{CMakePresets.json}{39}), so instrumented builds
never disturb the normal build directory.

All non-Qt dependencies arrive via FetchContent at exact tags, enumerated in
\cref{tab:overview-deps}; Qt6 is the only system dependency (\src{gui/CMakeLists.txt}{13}),
which makes the Qt5/Android fallback branch at \src{gui/CMakeLists.txt}{66} unreachable dead
code --- a safe deletion candidate.

\begin{designrule}[Hermetic dependency closure]
curl is built as a static, \code{HTTP\_ONLY} library (\src{CMakeLists.txt}{46}) with every
optional third-party backend --- libssh2, nghttp2, libidn2, libpsl, brotli, zstd --- forced
off. The in-tree comment records why (\src{CMakeLists.txt}{48}): several of these knobs
default to ON/AUTO, so curl links whatever it auto-detects on the build host, and
auto-detected Homebrew libraries on the macOS CI runner previously broke the final link.
The forced-off block converts an environment-dependent link into a reproducible one.
\end{designrule}

One knob in that block is inert: \cppinline|set(SSL_ENABLED ON)| at \src{CMakeLists.txt}{47}
is not a curl cache option at all --- the real switch, \code{CURL\_ENABLE\_SSL}, defaults ON
in the fetched source (\src{build/\_deps/curl-src/CMakeLists.txt}{725}). HTTPS works, but
through curl's defaults, not through line 48; the line misleads a reader about how TLS is
enabled.

Warning hygiene is strict tree-wide: \cppinline|-Wall -Wextra -Wpedantic -Werror| on
GCC/Clang (\src{CMakeLists.txt}{91}), \code{/W4 /WX} on MSVC (\src{CMakeLists.txt}{82}),
applied after the FetchContent blocks so vendored code is exempt --- with one deliberate
exception: the vendored \srcf{gui/qcustomplot.cpp} is compiled \code{-w}
(\src{gui/CMakeLists.txt}{20}) so \code{-Werror} can stay on for project code. Against this, a
comment in the same block claims ``We also run ASAN and clang-tidy for additional security''
(\src{CMakeLists.txt}{87}); nothing backs it --- no \code{.clang-tidy} file, no
\code{-fsanitize} flag, no lint preset, no CI job. Sanitizers and static analysis are
\fabsent, and the comment is stale. The layered link graph the build enforces
(\code{utils} $\rightarrow$ \code{sim} $\rightarrow$ \code{model} $\rightarrow$ the
\code{qtrocket\_core} static library at \src{core/CMakeLists.txt}{8}, shared by every
executable) is covered in \cref{sec:overview}; the bottom-layer invariant is stated in the
build itself (\src{core/utils/CMakeLists.txt}{16}) and holds --- nothing in \code{utils/} includes
a \code{model/} or \code{sim/} header.

\subsection{The test suite: six binaries, 247 cases}\label{sec:testing-suite}

Testing is GoogleTest throughout: six binaries totalling 247 \code{TEST}/\code{TEST\_F}
cases, summarized in \cref{tab:testing-binaries}.

\begin{table}[htbp]
\centering\small
\begin{tabularx}{\linewidth}{@{}l r X@{}}
\toprule
Binary & Cases & Focus \\
\midrule
\code{model\_tests} & 144 & Composite-tree mechanics (dirty flags, parallel-axis shifts, deep clone); part geometry against closed-form, numeric, and mesh oracles; Barrowman aero composition; motor burn; thrust-curve interpolation; JSON/RASP motor parsing; placement resolver and diagnostics \\
\code{sim\_tests} & 27 & US Standard Atmosphere tables; pluggable environment selection; RKF45 against closed-form dynamics; spherical gravity and geoid; integrator-selection guards \\
\code{integration\_tests} & 45 & End-to-end physics through the controller (11); motor-database persistence (9); \code{.qrd} persistence (14); placement invariance against a committed baseline, including the gated baseline regenerator (9); poke-through regressions (2) \\
\code{propagator\_tests} & 6 & Termination guards driven by a mock \code{Propagatable}: no-liftoff, non-finite state, max-sim-time backstop, clean abort on integrator exception \\
\code{cli\_tests} & 6 & Scripted REPL build--save--reload--fly; \code{key=value} and link-token validation; seat-vocabulary round trip; \code{checkdesign} overlap/air-gap diagnostics; RASP motor loading \\
\code{design\_matrix\_tests} & 19 & 25-fixture corpus; mass scaling laws; 1/4A$\rightarrow$M flight ladder (apogee monotonic in impulse); drag and integrator cross-checks; CLI-script rebuild of the committed corpus (3) \\
\midrule
total & 247 & \\
\bottomrule
\end{tabularx}
\caption{The six GoogleTest binaries and their coverage map. Case counts by grep of
\code{TEST}/\code{TEST\_F} at the pinned commit.}
\label{tab:testing-binaries}
\end{table}

The suite's character is the project's strongest engineering asset: physics is validated
against independent oracles --- closed-form inertia formulas, numeric disk integration, mesh
cross-checks (\cref{sec:model}) --- rather than recorded outputs; persistence is locked by
round-trips plus a committed bit-identical invariance baseline over the fixture corpus
(\cref{sec:persistence}); and each of the propagator's historical failure modes has a
dedicated regression (\cref{sec:sim}). Motor data is tested at three altitudes: the
thrustcurve.org JSON parsers are pure functions fed canned JSON string literals, no network
involved (\src{model/tests/ThrustCurveParserTests.cpp}{1}); the database's online-source
merge is driven through a \code{FakeThrustCurveAPI} implementing the
\code{model::ThrustCurveAPI} seam (\src{tests/MotorDatabasePersistenceTests.cpp}{106}); and
the design-matrix suite flies real motors end to end. Fixture location is wired by
absolute-path compile definitions (\code{QTROCKET\_DATA\_DIR},
\code{QTROCKET\_TEST\_DATA\_DIR}, \src{tests/CMakeLists.txt}{15}), so every binary runs from
any working directory.

Registration with ctest follows one pattern everywhere: each binary is registered per-case
via \code{gtest\_discover\_tests} and again as a single aggregate \code{qtrocket\_*} entry
(\src{core/model/tests/CMakeLists.txt}{29}; \src{tests/CMakeLists.txt}{25}), giving seven
aggregate entries in total. The design-matrix binary carries the seventh --- a third
registration of the same executable, \code{qtrocket\_design\_matrix\_heavy\_tests}, which
reruns only the \code{FlightMatrix}/\code{Atmosphere} sweeps under the ctest label
\code{heavy} with \code{QTROCKET\_FULL\_LADDER=1} in its environment
(\src{tests/CMakeLists.txt}{97}). That environment variable is the whole mechanism: the
flight sweeps ask \code{motorsToFly()} which motors to fly, and it returns either one motor
per impulse class or the class's complete ladder (\src{tests/DesignMatrixTests.cpp}{376}):

\begin{minted}{cpp}
std::vector<std::string> motorsToFly(const Ladder& lad)
{
   if(std::getenv("QTROCKET_FULL_LADDER") != nullptr)
      return lad.motors;
   return { lad.motors.front() };
}
\end{minted}

\begin{keyinsight}[One binary, three ctest entries]
The double registration means an unfiltered \code{ctest} runs every case twice --- once
discovered, once inside its aggregate. The documented workflows always filter
(\code{-R 'qtrocket\_*'} for everything, \code{-LE heavy} for the fast loop, \code{-L heavy}
for the sweeps alone, \src{tests/CMakeLists.txt}{86}), so this is a footgun rather than a
bug --- but one a newcomer running bare \code{ctest} will step on. The heavy/fast split
itself is well designed: one binary, an env-var gate, and a ctest label buy a fast default
loop and an on-demand full-ladder sweep with no code duplication.
\end{keyinsight}

\subsection{Coverage tooling}\label{sec:testing-coverage}

Coverage is Clang/LLVM source-based and gated hard: \code{QTROCKET\_ENABLE\_COVERAGE}
fatally errors on any non-Clang compiler (\src{CMakeLists.txt}{95}), then applies
\cppinline|-O0 -g -fprofile-instr-generate -fcoverage-mapping| (\src{CMakeLists.txt}{102}).
A custom \code{coverage} target shells out to \srcf{cmake/RunLlvmCoverage.cmake}, which runs
ctest under \code{LLVM\_PROFILE\_FILE} (\src{cmake/RunLlvmCoverage.cmake}{33}), aborts on
any test failure (\src{cmake/RunLlvmCoverage.cmake}{39}), merges profiles with
\code{llvm-profdata}, and emits both a text summary and a branch-annotated, demangled HTML
report (\src{cmake/RunLlvmCoverage.cmake}{99}). This is genuinely good tooling --- and no CI
job runs it, so coverage numbers exist only on a developer's machine. Verdict: \fpartial,
for that reason and for a real defect in its binary list.

\begin{hardfail}[Coverage target out of sync with the design-matrix suite]
\code{QTROCKET\_COVERAGE\_BINARIES} (\src{CMakeLists.txt}{139}) and the target's
\code{DEPENDS} (\src{CMakeLists.txt}{158}) name five test binaries ---
\code{design\_matrix\_tests} is missing from both --- while the ctest regex the script is
handed, \cppinline|^qtrocket_.*tests$| (\src{CMakeLists.txt}{152}), matches both
design-matrix entries. On a fresh \code{build-coverage} tree, \code{--target coverage}
therefore builds only the five listed binaries, ctest then tries to run the design-matrix
entries whose executable was never built, and the script aborts with \code{FATAL\_ERROR}
(\src{cmake/RunLlvmCoverage.cmake}{39}). Even when the binary happens to exist, the heavy
full-ladder sweep runs under instrumentation, and functions unique to that suite are never
attributed in the report. \fail
\end{hardfail}

\subsection{Continuous integration}\label{sec:testing-ci}

CI is a GitHub Actions workflow triggered on push and pull request to \code{development}
only (\src{.github/workflows/cmake-multi-platform.yml}{13}). A four-job matrix with
fail-fast off --- Linux/GCC, Linux/Clang, macOS/AppleClang, and Windows/MSVC, each driving
its release preset (\src{.github/workflows/cmake-multi-platform.yml}{29}) --- installs Qt
6.10.2 (\src{.github/workflows/cmake-multi-platform.yml}{45}), caches the FetchContent
dependencies via ccache, builds the entire tree including the GUI and visualizer
(\src{.github/workflows/cmake-multi-platform.yml}{68}), and runs
\code{ctest -R 'qtrocket\_*'} (\src{.github/workflows/cmake-multi-platform.yml}{74}). A
fifth job boots a FreeBSD 15.0 QEMU virtual machine on the Ubuntu runner
(\src{.github/workflows/cmake-multi-platform.yml}{94}), installs Qt from \code{pkg}
(\src{.github/workflows/cmake-multi-platform.yml}{102}), and runs the same preset and
filter. Exercising GCC, Clang, AppleClang, and MSVC under \code{-Werror}/\code{/WX} on
every merge request is real rigor for a project of this size, and the FreeBSD job goes
beyond what most C++ projects of any size attempt.

Two blemishes. First, the ctest regex matches the \code{heavy} entry, so all five jobs ---
the QEMU VM included --- fly every motor of every impulse-class ladder on every pull
request: exactly the cost the \code{-LE heavy} fast loop exists to avoid. Second, the test
step's comment still names only four suites, ``(model, sim, integration, propagator)''
(\src{.github/workflows/cmake-multi-platform.yml}{72}), predating the cli and design-matrix
binaries. Beyond running tests, CI publishes nothing: no coverage upload, no sanitizer or
lint job, no artifacts, no releases.

\subsection{The utils layer}\label{sec:testing-utils}

The bottom layer is deliberately small --- four components, Qt-free --- and its quality is
more mixed than the layers above it.

\begin{description}
\item[\code{Logger}] A Meyers-singleton logger (\src{utils/Logger.cpp}{19}) with five
cumulative levels, an intentional-fallthrough switch under a mutex
(\src{utils/Logger.cpp}{35}), and dual sinks (stdout plus \code{log.txt}). Three defects:
\code{setLogLevel} writes the level without taking the mutex \code{log()} holds while
reading it (\src{utils/Logger.cpp}{77}) --- benign today, since every call site is startup
or test setup, but a latent data race; \code{log.txt} opens relative to the process working
directory (\src{utils/Logger.cpp}{25}), so the file lands wherever the app was launched
from; and entries carry level tags but no timestamps.
\item[\code{Bin}] A lower-bound interval map (\src{utils/Bin.h}{15}) whose sole consumers
are the four static US Standard Atmosphere layer tables
(\src{sim/USStandardAtmosphere.h}{28}, initialized at
\src{sim/USStandardAtmosphere.cpp}{78}). It throws \code{std::out\_of\_range} below the
lowest bin, but both \code{operator[]} and \code{getBinBase} dereference
\cppinline|bins.begin()| before any empty check (\src{utils/Bin.cpp}{50}) --- undefined
behavior on an empty \code{Bin} --- and \code{insert} re-sorts the entire vector on every
call (\src{utils/Bin.cpp}{44}). Both limitations are acknowledged in-tree
(\src{utils/Bin.h}{17}) and unreachable today, the only consumers being statically
populated; neither is tested.
\item[math headers] Global-namespace Eigen aliases (\src{utils/math/MathTypes.h}{16}), the
physical-constants header whose \code{earthGM} feeds the spherical gravity model
(\src{utils/math/Constants.h}{26}), and a ULP-scaled \code{floatingPointEqual}
(\src{utils/math/UtilityMathFunctions.h}{21}). Dead weight: a 47-line commented-out block
in \srcf{utils/math/MathTypes.h} (\src{utils/math/MathTypes.h}{24}), and
\code{earthGM\_km} (\src{utils/math/Constants.h}{27}), consumerless since the gravity path
went pure SI.
\item[\code{CurlConnection}] An RAII libcurl GET wrapper: copying deleted because it would
double-free the handle (\src{utils/CurlConnection.h}{23}), reference-counted global init per
instance (\src{utils/CurlConnection.cpp}{29}), and hard $10\unit{s}$ connect /
$30\unit{s}$ total timeouts with \code{NOSIGNAL}, justified in-tree by the synchronous
GUI-thread call path (\src{utils/CurlConnection.cpp}{54}). Transport errors and non-2xx
statuses map to a logged error and an empty return (\src{utils/CurlConnection.cpp}{68});
TLS verification is never set explicitly, relying on libcurl's default-on peer and host
checks.
\end{description}

\subsection{Untested surfaces and dormant machinery}\label{sec:testing-gaps}

The coverage map has three empty regions. \code{utils/} has no test subdirectory at all
(\src{core/utils/CMakeLists.txt}{1}): Logger's level semantics, Bin's empty-container edge, and
\code{floatingPointEqual} are directly untested, Bin being exercised only through sim's
atmosphere tests (\src{sim/tests/USStandardAtmosphereTests.cpp}{5}). \code{CurlConnection}
is never executed by any test --- the network is stubbed one layer up, at the
\code{ThrustCurveAPI} seam (\cref{sec:motors}) --- so the one component that talks to the
outside world is validated only in production. And the entire Qt surface, \code{gui/} and
\code{visualizer/}, has zero automated tests: CI compiles both on every platform
(\src{.github/workflows/cmake-multi-platform.yml}{68}) but runs only the ctest suites
(\src{.github/workflows/cmake-multi-platform.yml}{74}), so the GUI is built five ways and
exercised zero (\cref{sec:interfaces}). There is likewise no fuzzing of the XML, JSON, or
RASP loaders that parse untrusted files.

Distribution is where the unconsumed-machinery motif is starkest. The sole install rule
covers the GUI target only (\src{gui/CMakeLists.txt}{96}), referencing a
\code{CMAKE\_INSTALL\_LIBDIR} that is never defined because \code{GNUInstallDirs} is never
included; \code{qtrocket-cli}, \code{qtrocket-visualizer}, and the bundled \code{data/}
motor files are not installed, there is no CPack, no desktop file, no installer of any
kind, and the macOS bundle identifier is the Qt-template placeholder
\code{my.example.com} (\src{gui/CMakeLists.txt}{89}). The project version, 0.1, exists only in
\cppinline|project()| (\src{CMakeLists.txt}{3}) and the bundle metadata --- the About dialog does not surface it and no \code{--version} flag
or \code{--version} flag surfaces it. Internationalization is vestigial in the same way:
\code{qt\_create\_translation} generates \code{QM\_FILES} from the template
\code{qtrocket\_en\_US.ts} (\src{gui/CMakeLists.txt}{65}, \src{gui/CMakeLists.txt}{18}), but
nothing references \code{QM\_FILES}, so no translation is ever built or shipped ---
\code{LinguistTools} is a required find-package component (\src{gui/CMakeLists.txt}{12}) in
service of an inert call. Where OpenRocket ships platform installers, a bundled runtime, and
roughly fifteen languages, QtRocket currently ships nothing at all; the comparison is drawn
in \cref{sec:comparison}, and the consolidated remediation list in \cref{sec:roadmap}.
